% Chapter 21: The Erd\H{o}s Discrepancy Problem
\let\LocalMacro\relax

\chapter{Erdős Discrepancy Problem}

\section{The Problem}

\subsection{Informal}
Imagine assigning either plus one or minus one to every counting number: 1, 2, 3, 4, and so on forever. Now look at every evenly-spaced pattern you can find---the even numbers, the multiples of 3, the multiples of 7, and so on---and add up the signs along each pattern. Can you assign the signs so that \emph{every} such sum stays small, no matter how far out you go?

\subsection{Formal}
The Erd\H{o}s discrepancy problem asks:
\begin{conjecture}[Erd\H{o}s, 1957]
For every completely multiplicative function $f: \mathbb{Z} \to \{-1, +1\}$, there exist integers $x$ and $d$ such that
\[
\left| \sum_{k=1}^{x} f(kd) \right| > C
\]
for any fixed constant $C$. In other words, the partial sums along arithmetic progressions cannot be uniformly bounded.
\end{conjecture}

\subsection{Historical Context}
Paul Erdős posed this question in 1957 and offered \$500 for a solution. The problem resisted for 58 years---it is one of the simplest-to-state, hardest-to-solve puzzles in all of mathematics. Erdős called it one of his favorites.

\subsection{What Was Known}
For decades, mathematicians proved results for special cases---bounded discrepancy for specific sequences, or discrepancy bounds for restricted sets of progressions. The full conjecture, that no assignment of plus and minus can keep all partial sums bounded, remained completely open for 58 years.

\subsection{Why It Matters}
The discrepancy problem reveals deep truths about the structure of the integers and the limits of balance. It connects to the theory of random sequences, the distribution of primes, and the foundations of probability theory. Tao's proof introduced new techniques in additive combinatorics that have been applied to many other problems.

\subsection{Simplified Version}
For \emph{periodic} completely multiplicative sequences, the discrepancy is obviously infinite: the sequence repeats, so the sums over arithmetic progressions grow without bound. The hard part was proving that \emph{every} completely multiplicative sequence (including non-periodic ones) has unbounded discrepancy. This was proved by Terence Tao in 2015 using a combination of Fourier analysis and the probabilistic method.

\section{Mathematical Theory}


\subsection{Prerequisites}
This chapter assumes familiarity with:
\begin{itemize}
\item Basic analysis (sequences, series, $\limsup$)
\item Elementary combinatorics
\item Fourier analysis on finite abelian groups
\item Basic probability theory
\end{itemize}

\subsection{Discrepancy and Multiplicative Functions}
A sequence $x: \mathbb{N} \to \{-1, +1\}$ is \emph{completely multiplicative} if $x(mn) = x(m)x(n)$. For such sequences, the discrepancy simplifies:
\[
\sum_{k=1}^N x_{kd} = x_d \sum_{k=1}^N x_k
\]
so the discrepancy is controlled by the partial sums of the original sequence.

\begin{conjecture}[Erdős Discrepancy Conjecture]
For any sequence $x_n \in \{-1, +1\}$, the discrepancy $D = \infty$.
\end{conjecture}

\subsection{Connection to the Liouville Function}
The Liouville function $\lambda(n) = (-1)^{\Omega(n)}$ (where $\Omega(n)$ counts prime factors with multiplicity) is completely multiplicative. The Erdős discrepancy conjecture for $\lambda$ is closely related to Chowla's conjecture on correlations of multiplicative functions, but the general conjecture is much stronger and applies to all completely multiplicative sequences, not just the Liouville function.

\subsection{Previous Approaches}
\begin{itemize}
\item \textbf{Beck-Fiala theorem \cite{beckfiala1981}}: For set systems with bounded degree, discrepancy $O(\sqrt{d} \log n)$.
\item \textbf{Spencer's theorem \cite{spencer1985}}: Six standard deviations suffice ($O(\sqrt{n})$).
\item \textbf{Banaszczyk's theorem (1998)}: Discrepancy $O(\sqrt{\log n})$ for arbitrary vectors.
\item \textbf{Computer searches}: Long sequences with low discrepancy were found (length 1124 with discrepancy 2), suggesting the conjecture might be false.
\end{itemize}

\section{The Solution}

\subsection{The Big Idea}
Imagine shuffling a deck of cards forever---plus one, minus one, plus one, minus one, assigned to every counting number. Can you keep \emph{every} evenly-spaced sum (even numbers, multiples of 3, multiples of 7, \ldots) small for all time? Tao showed this is impossible by a beautiful indirect argument: if you \emph{could} keep all sums bounded, then zooming out and averaging the sequence over larger and larger blocks of primes would make the sequence more and more predictable---its ``information content'' would drop. It would be forced to look like a periodic pattern based on prime numbers (a Dirichlet character). But periodic patterns always have unbounded sums along their own period. You cannot hide from that. It is like proving that no matter how you shuffle a deck of cards, if you zoom out enough, patterns \emph{must} emerge, and those patterns inevitably blow up the discrepancy.

\subsection{What Was Stopping Progress}
The problem was open for 58 years because the tools of combinatorics---Beck--Fiala, Spencer's ``six standard deviations,'' Banaszczyk's bounds---work for \emph{set systems} with bounded degree. The Erd\H{o}s problem has \emph{all arithmetic progressions} as its set system, which has unbounded degree. These methods were fundamentally the wrong shape for the problem. Fourier analysis could detect periodicity but could not handle non-periodic sequences. Purely combinatorial arguments lacked the leverage to force \emph{global} unboundedness from \emph{local} constraints. It was like trying to prove a rope is infinitely long by only examining small sections of it.

\subsection{The Breakthrough}
Tao's key idea was the \emph{entropy decrement argument}: an averaging technique that shows bounded discrepancy forces the sequence to become more and more predictable as you average over larger primes. Once the sequence is predictable enough, it must correlate with a Dirichlet character, and that character has unbounded discrepancy---a contradiction. The proof also required proving a logarithmically averaged version of Elliott's conjecture (a far-reaching generalization of Chowla's conjecture), which formalized the statement that correlation with a Dirichlet character propagates.

\subsection{How It Works}

Before Tao, the Erd\H{o}s discrepancy conjecture had been open for 58 years with virtually no progress on the general case. Partial results existed for special sequences---for instance, it was known that periodic completely multiplicative sequences have unbounded discrepancy, and that certain structured sequences (those correlated with Dirichlet characters) cannot maintain low discrepancy. But these special cases shed no light on the core difficulty: proving that \emph{every} completely multiplicative $\pm 1$ sequence, no matter how cleverly constructed, must eventually exhibit large partial sums along some arithmetic progression.

Classical tools in combinatorics---the Beck--Fiala theorem, Spencer's ``six standard deviations suffice'' result, and Banaszczyk's bounds---address discrepancy for \emph{set systems} (bounded-degree hypergraphs) and produce bounds of the form $O(\sqrt{n})$ or $O(\sqrt{\log n})$. These methods are fundamentally ill-suited to the Erd\H{o}s problem because the ``set system'' here consists of \emph{all} arithmetic progressions, which have unbounded degree. Fourier analysis on the integers, while useful for detecting periodicity, cannot directly handle non-periodic completely multiplicative sequences. Purely combinatorial arguments lacked the leverage to force global unboundedness from local constraints.

Tao's 2015 proof introduced a new probabilistic--analytic technique:
\begin{enumerate}
\item \textbf{The entropy decrement argument:} This is a new averaging technique. If the discrepancy were bounded, one could average the sequence over large primes and show that the ``information content'' (entropy) of the sequence decreases. This forces the sequence to correlate with a Dirichlet character of bounded conductor.
\item \textbf{Reduction to Elliott's conjecture:} Tao proved a logarithmically averaged version of Elliott's conjecture (a far-reaching generalization of Chowla's conjecture on correlations of multiplicative functions), which formalized the statement that correlation with a Dirichlet character propagates.
\item \textbf{Forcing unbounded discrepancy:} Once correlation with a Dirichlet character is established, the discrepancy is forced to be unbounded because Dirichlet characters have unbounded partial sums along certain arithmetic progressions (by Dirichlet's theorem on primes in arithmetic progressions).
\end{enumerate}

Suppose you could assign $+1$ and $-1$ to every counting number in a way that keeps \emph{every} evenly-spaced sum small forever. Tao showed this is impossible by an indirect argument. First, he showed that such a carefully balanced sequence would have to ``look like'' a periodic pattern based on prime numbers (a Dirichlet character). But periodic patterns always have unbounded sums along their own period---you can't hide from that. The key new idea is the \emph{entropy decrement}: imagine zooming out and averaging the sequence over larger and larger blocks of primes. If the discrepancy were bounded, this averaging would make the sequence more and more predictable (its entropy would drop), forcing it to look like a Dirichlet character---and that leads to a contradiction.

\bigskip

Terence Tao announced a proof on September 18, 2015, using a combination of:

\begin{enumerate}
\item \textbf{Entropy decrement argument}: A new averaging technique showing that if discrepancy is bounded, the sequence must correlate with a Dirichlet character.
\item \textbf{Elliott's conjecture}: A generalization of Chowla's conjecture on correlations of multiplicative functions, proved by Tao using entropy methods.
\item \textbf{Fourier analysis and nilsequences}: Showing that bounded discrepancy forces correlation with a nilsequence, which then forces unbounded discrepancy via Gowers norms.
\end{enumerate}

\begin{theorem}[Tao \cite{tao2016discrepancy}]
For any sequence $x_n \in \{-1, +1\}$, the discrepancy is infinite.
\end{theorem}

\begin{proof}[Sketch]
Assume for contradiction that $D < \infty$. Then:
\begin{enumerate}
\item By the entropy decrement argument, $x$ correlates with a Dirichlet character $\chi$ of bounded conductor.
\item By Elliott's conjecture (proved via entropy), this correlation forces $x$ to behave like $\chi$.
\item But a completely multiplicative function of bounded discrepancy must be close to a Dirichlet character, and such characters have unbounded discrepancy along some arithmetic progression (by Dirichlet's theorem on primes in AP).
\item Contradiction.
\end{enumerate}
\end{proof}

\begin{highlightbox}
Tao's proof was remarkable because it combined probabilistic combinatorics (entropy), analytic number theory (Elliott's conjecture), and ergodic theory (nilsequences). The proof is non-constructive and does not give explicit bounds on how fast the discrepancy grows.
\end{highlightbox}

\subsection{Subsequent Developments}
\begin{itemize}
\item \textbf{Polymath project}: Verified and streamlined the proof.
\item \textbf{Quantitative bounds}: No explicit growth rate for discrepancy is known (the proof gives a tower-of-exponentials type bound).
\item \textbf{Variants}: Erdős discrepancy for complex units, for sequences in $\{-1,0,1\}$, etc.
\end{itemize}

\subsection{After the Proof}

The Erd\H{o}s discrepancy problem is completely settled: every completely multiplicative $\pm 1$ sequence has unbounded discrepancy. The conjecture is a theorem, and there is nothing left to prove about the qualitative question. The Polymath project verified and streamlined Tao's proof, making it accessible to a broader community.

One important limitation is that the proof is non-constructive. It shows discrepancy must be infinite, but it does not say how fast it grows. The best bound coming from the proof is tower-of-exponentials: an absurdly fast-growing function that is far from what computer searches suggest. Computer experiments find sequences of length over 1000 with discrepancy 2, hinting that the true growth rate might be much slower. Determining the exact growth rate of the discrepancy is a major open problem.

Several natural extensions remain open. What happens if you allow the sequence values to be complex numbers on the unit circle instead of just $\pm 1$? What about sequences valued in $\{-1, 0, 1\}$? These variants have partial results but are not fully resolved. The entropy decrement technique that Tao introduced has also been applied to other problems in additive combinatorics---particularly to variants of Chowla's conjecture on correlations of the Liouville function---but the full range of consequences is still being explored.

\section{Impact}
\begin{itemize}
\item \textbf{Combinatorics}: First major application of entropy methods to discrepancy theory.
\item \textbf{Number theory}: Elliott's conjecture (proved en route) is a major result on correlations of multiplicative functions.
\item \textbf{Computer science}: Implications for online algorithms and streaming algorithms.
\item \textbf{Erdős legacy}: Erdős's \$500 prize was finally claimed (Tao donated it to charity).
\end{itemize}

\section{Computational Verification}

The Erd\H{o}s discrepancy problem can be explored computationally by searching for low-discrepancy sequences and verifying that discrepancy eventually grows.

\subsection{Computing Discrepancy}

\begin{verbatim}
def discrepancy(sequence, max_d=None):
    """
    Compute the discrepancy of a completely multiplicative sequence.
    D = max over d, x of |sum_{k=1}^{x} f(k*d)|
    """
    N = len(sequence)
    if max_d is None:
        max_d = N
    D = 0
    for d in range(1, max_d + 1):
        s = 0
        for k in range(1, N // d + 1):
            s += sequence[k * d - 1]
            D = max(D, abs(s))
    return D

# Example: constant sequence f(n) = 1
seq = [1] * 1000
print(f"Constant sequence, N=1000: discrepancy = {discrepancy(seq, 50)}")
# Should be large (unbounded)
\end{verbatim}

\subsection{Searching for Low-Discrepancy Sequences}

Computer searches found sequences of length 1124 with discrepancy 2, suggesting the growth rate might be slow.

\begin{verbatim}
import random

def find_low_discrepancy(N, target_D=2, attempts=100000):
    """Random search for completely multiplicative sequences with low discrepancy."""
    best_seq = None
    best_D = float('inf')

    for _ in range(attempts):
        # Random completely multiplicative: assign random +/-1 to primes
        seq = [0] * N
        primes = sieve_of_eratosthenes(N)
        prime_vals = {p: random.choice([-1, 1]) for p in primes}
        for i in range(1, N + 1):
            val = 1
            n = i
            for p in primes:
                if p * p > n:
                    break
                while n % p == 0:
                    val *= prime_vals[p]
                    n //= p
            if n > 1 and n in prime_vals:
                val *= prime_vals[n]
            seq[i - 1] = val

        D = discrepancy(seq, min(50, N))
        if D < best_D:
            best_D = D
            best_seq = seq[:]

    return best_seq, best_D
\end{verbatim}

\subsection{Liouville Function Discrepancy}

The Liouville function $\lambda(n) = (-1)^{\Omega(n)}$ is a natural test case.

\begin{verbatim}
def liouville(n):
    """Compute the Liouville function lambda(n)."""
    if n == 1:
        return 1
    count = 0
    d = 2
    while d * d <= n:
        while n % d == 0:
            count += 1
            n //= d
        d += 1
    if n > 1:
        count += 1
    return (-1) ** count

N = 500
seq = [liouville(n) for n in range(1, N + 1)]
D = discrepancy(seq, 100)
print(f"Liouville function, N={N}, discrepancy (d<=100) = {D}")
\end{verbatim}

\begin{highlightbox}
Students should run these scripts to observe: (1)~the Liouville function has slowly growing discrepancy, (2)~random completely multiplicative sequences tend to have low discrepancy for small $N$, and (3)~the growth rate of discrepancy is much slower than the tower-of-exponentials bound from Tao's proof.
\end{highlightbox}

\section{Exercises}

\begin{exercise}[Basic]
Show that the sequence $x_n = 1$ for all $n$ has discrepancy $D = \infty$.
\end{exercise}

\begin{exercise}[Basic]
Show that any periodic sequence $x_n \in \{-1, +1\}$ has infinite discrepancy.
\end{exercise}

\begin{exercise}[Intermediate]
Let $x_n$ be completely multiplicative with $x_p = 1$ for all primes $p \equiv 1 \pmod{4}$ and $x_p = -1$ for $p \equiv 3 \pmod{4}$. Show that $x_n$ has infinite discrepancy.
\end{exercise}

\begin{exercise}[Challenge]
Explain the entropy decrement argument in your own words: how does averaging over large primes reduce the entropy of the sequence?
\end{exercise}

\section{Timeline}

\begin{tabularx}{\linewidth}{lX}
\toprule
\textbf{Year} & \textbf{Event} \\ \midrule
1957 & Erdős poses the discrepancy problem \\ 
1950s & Erdős offers \$500 prize \\ 
1981 & Beck-Fiala theorem \\ 
1985 & Spencer's "six standard deviations" theorem \\ 
2015 & Tao proves the Erdős discrepancy conjecture \\ 
2016 & Polymath project streamlines the proof \\ \bottomrule
\end{tabularx}

\section{Citations}

\begin{enumerate}
\item Tao, T. (2016). ``The Erdős discrepancy problem.'' Discrete Analysis, 1, 29 pp.
\item Tao, T. (2016). ``The logarithmically averaged Chowla and Elliott conjectures for two-point correlations.'' Forum of Mathematics, Pi, 4, e8.
\item Erdős, P. (1957). ``Some unsolved problems.'' Michigan Mathematical Journal, 4(3), 291--300.
\item Beck, J., \& Fiala, T. (1981). ``Integer-making theorems.'' Discrete Applied Mathematics, 3(1), 1--8.
\item Spencer, J. (1985). ``Six standard deviations suffice.'' Transactions of the AMS, 289(2), 679--706.
\end{enumerate}
